\QHeading{第4题　计算填表。}
\par\noindent


\DataTable{4}{出行交通工具 & 速度 & 行驶时间 & 行驶总路程\\\hline
二轮电动车 & 300米/分 & 25分钟 & {\color{answerred}7500米}\\\hline
家用轿车 & 90千米/时 & {\color{answerred}5小时} & 450千米\\\hline
大巴客车 & {\color{answerred}66千米/时} & 6小时 & 396千米\\\hline
山地车 & 220米/分 & 40分钟 & {\color{answerred}8800米}\\\hline}
\par

\par\noindent

\Answer{
\QHeading{计算过程：第4题：计算填表}
\begin{center}
\begin{adjustbox}{max width=\linewidth}\begin{tabular}{c|c|c|c}
\hline
出行交通工具&速度&行驶时间&行驶总路程\\\hline
二轮电动车&300米/分&25分钟&7500米\\
家用轿车&90千米/时&5小时&450千米\\
大巴客车&66千米/时&6小时&396千米\\
山地车&220米/分&40分钟&8800米\\\hline
\end{tabular}\end{adjustbox}
\end{center}
\begin{enumerate}
\item \(\begin{aligned}[t]&\phantom{={}}300\times25\\&=300\times(20+5)\\&=300\times20+300\times5\\&=6000+1500\\&=7500\text{（米）}\text{。}\end{aligned}\)
\item \(\begin{aligned}[t]&\phantom{={}}450\div90\\&=(450\div10)\div(90\div10)\\&=45\div9\\&=5\text{（小时）}\text{。}\end{aligned}\)
\item \(\begin{aligned}[t]&\phantom{={}}396\div6\\&=(360+36)\div6\\&=360\div6+36\div6\\&=60+6\\&=66\text{（千米/时）}\text{。}\end{aligned}\)
\item \(\begin{aligned}[t]&\phantom{={}}220\times40\\&=(22\times10)\times(4\times10)\\&=(22\times4)\times(10\times10)\\&=88\times100\\&=8800\text{（米）}\text{。}\end{aligned}\)
\end{enumerate}

}

\QHeading{第5题　解决问题。}

\Problem{①小宇和小琪两人同时从小区一段道路的两端相向而行，小宇每分钟走70米，小琪每分钟走65米，8分钟后两人相遇。小区这段道路全长多少米？}
\par


\MethodTitle{方法一}
\Answer{\Formula{\begin{aligned}[t]&\phantom{={}}(70+65)\times8\\&=135\times8\\&=1080\text{（米）}\end{aligned}}\par 答：这段道路全长1080米。}

\MethodTitle{方法二：分别计算两部分，再求和}
\Answer{不先合并每份的数量，而是分别求出两部分的总数，再相加。}
\Answer{\Formula{\begin{aligned}[t]&\phantom{={}}70\times 8+65\times 8\\
&=560+520\\
&=1080\text{（米）}\end{aligned}}}
\Answer{答：这段道路全长1080米。}
\Problem{②超市采购一批物资。每包湿巾24元，采购了25包，采购洗手液花费了360元。采购湿巾比采购洗手液多花多少钱？}

\MethodTitle{方法一}
\Answer{\Formula{\begin{aligned}[t]&\phantom{={}}24\times25-360\\&=600-360\\&=240\text{（元）}\end{aligned}}\par 答：采购湿巾比采购洗手液多花240元。}

\MethodTitle{方法二：先把360元折算成湿巾的包数}
\Answer{360元相当于15包湿巾的费用，采购25包就多10包的价钱。}
\Answer{\Formula{\begin{aligned}[t]&\phantom{={}}24\times (25-360\div 24)\\
&=24\times (25-15)\\
&=24\times 10\\
&=240\text{（元）}\end{aligned}}}
\Answer{答：采购湿巾比采购洗手液多花240元。}
\Problem{③两辆汽车从相距500千米的两地同时相向开出，甲车的速度为60千米/时，乙车的速度为40千米/时。几小时后两车相遇？}

\MethodTitle{方法一}
\Answer{\Formula{\begin{aligned}[t]&\phantom{={}}500\div(60+40)\\&=500\div100\\&=5\text{（小时）}\end{aligned}}\par 答：5小时后两车相遇。}

\MethodTitle{方法二：列表累加路程}
\Answer{每过1小时，甲车走60千米、乙车走40千米，合起来走100千米。把两车的行程逐时相加。}
\Answer{\textbf{每小时合计行程}\par \Formula{\begin{aligned}[t]&\phantom{={}}60+40\\
&=100\text{（千米）}\end{aligned}}}
\Answer{\textbf{5小时合计行程}\par \Formula{\begin{aligned}[t]&\phantom{={}}100\times 5\\
&=500\text{（千米）}\end{aligned}}}
\Answer{1、2、3、4、5小时，两车合计行驶100、200、300、400、500千米；第5小时刚好走完两地之间的路程。}
\Answer{答：5小时后两车相遇。}

